(1) Euler–Lagrange变分给 $\boxed{\mu y_{tt}+EI y_{xxxx}=0}$。它是空间四阶、时间二阶的线性波动方程，弯曲刚度而不是一维横向绳张力提供回复力。对任意满足适当边界条件的场，总弹性能非负。

(2) 平面波 $e^{i(kx-\omega t)}$ 使 $\mu\omega^2=EI k^4$，正频分支 $\omega=\sqrt{EI/\mu}\,k^2$。所以 $v_p=\omega/k=\sqrt{EI/\mu}\,k$，$v_g=d\omega/dk=2\sqrt{EI/\mu}\,k=2v_p$，该理想模型在高波数出现越来越快的波包。

(3) 乘PDE以 $y_t$，整理为 $\partial_t[\mu y_t^2/2+EI y_{xx}^2/2]+\partial_x[EI(y_{xxx}y_t-y_{xx}y_{xt})]=0$。对 $y=A\cos(kx-\omega t)$，平均能密度 $\bar e=\mu A^2\omega^2/2$，能流 $\bar S=EI A^2k^3\omega$，故 $\bar S/\bar e=2EI k^3/(\mu\omega)=2\omega/k=v_g$。

(4) 简支边界为 $y(0)=y(L)=0$ 且弯矩 $EI y_{xx}(0)=EI y_{xx}(L)=0$。允许模态 $Y_n=\sin(n\pi x/L)$，频率 $\boxed{\omega_n=\sqrt{EI/\mu}(n\pi/L)^2}$。当波长与梁厚接近时，Euler–Bernoulli忽略的转动惯量、剪切变形和三维弹性使高速渐近不再可靠。
